\section{Contexto y motivación}
\label{sec:contexto_motivacion}

La capacidad de cómputo se ha convertido en un input económico relevante para el
desarrollo y despliegue de sistemas modernos de inteligencia artificial. El aumento
de los recursos computacionales disponibles ha acompañado la expansión del
\textit{deep learning} y de los modelos de gran escala
\parencite{sevilla2022compute}. Para las empresas que no mantienen infraestructura
propia, una parte de esta capacidad se adquiere mediante servicios de
\textit{cloud computing} y plataformas especializadas, donde el coste se expresa
habitualmente por GPU-hora. La literatura sobre mercados de cloud ha analizado
problemas de asignación, contratación y pricing dinámico de estos recursos
\parencite{xu2013dynamiccloud,li2016spotpricing,wu2019cloudpricing}.

Desde la perspectiva de gestión del riesgo, una empresa con necesidades recurrentes
de cómputo no está expuesta sólo al precio observado en una fecha concreta, sino al
coste medio pagado durante un horizonte. Esta característica motiva el uso de una
opción asiática aritmética, cuyo payoff depende del promedio de una secuencia de
precios. Para $N$ fechas de fixing $t_1,\ldots,t_N$, el contrato considerado es

\begin{equation}
    H_T
    =
    \left(
        \frac{1}{N}\sum_{i=1}^{N} P_{t_i}-K
    \right)^+,
\end{equation}

donde $P_t$ representa el precio por GPU-hora y $K$ el strike. El contrato se
interpreta como un instrumento hipotético de cobertura del coste medio de capacidad,
no como evidencia de que exista ya un mercado líquido y completo de derivados sobre
compute. En 2026 se anunciaron iniciativas dirigidas a desarrollar futuros
referenciados a índices de alquiler de GPU, pero la historia disponible todavía no
permite tratar este mercado como una estructura madura de derivados
\parencite{cme_compute_2026}.

El problema presenta además una dimensión computacional. La opción es dependiente
de trayectoria y su valoración de referencia se obtiene mediante simulación. Esto
hace natural estudiar redes neuronales como \textit{surrogates} de la función de
pricing y, en particular, si incorporar información diferencial durante el
entrenamiento mejora la aproximación de la sensibilidad relevante para cobertura.
Differential Machine Learning (DML) extiende el aprendizaje supervisado convencional
utilizando simultáneamente etiquetas de valor y de derivadas
\parencite{huge_savine_2020}. El presente TFM combina esta idea con un entorno de
valoración RQMC y una dinámica Log-Ornstein--Uhlenbeck para estudiar no sólo la
precisión numérica del precio y la Delta, sino también su utilidad económica en
cobertura.

\section{Problema de investigación}
\label{sec:problema_investigacion}

El trabajo combina dos preguntas. La primera es numérica: si un MLP convencional y
DML, comparados bajo la misma arquitectura y protocolo, aproximan de forma diferente
el precio y la Delta de una opción asiática. La segunda es financiera: si una mejora
en la estimación de la Delta se traduce efectivamente en una reducción del error de
cobertura.

El motor RQMC genera, para cada estado $\mathbf{x}$, una referencia de precio
$V(\mathbf{x})$ y una sensibilidad

\begin{equation}
    \Delta(\mathbf{x})
    =
    \frac{\partial V(\mathbf{x})}{\partial P}.
\end{equation}

El MLP convencional aprende a partir de pares $(\mathbf{x}_i,V_i)$, mientras que
DML utiliza además la etiqueta diferencial $\Delta_i$. La comparación se diseña
para aislar esta diferencia y evitar que una eventual ventaja se deba a arquitecturas
o procedimientos de selección distintos.

La \textbf{pregunta principal de investigación} es:

\begin{quote}
\emph{¿En qué medida la incorporación de información diferencial mediante
Differential Machine Learning mejora, respecto a un MLP convencional, la
aproximación del precio y especialmente de la Delta de una opción asiática
aritmética sobre capacidad de cómputo?}
\end{quote}

Como complemento, se plantea una segunda pregunta de naturaleza económica:

\begin{quote}
\emph{¿Se traduce una estimación más precisa de la Delta en una reducción efectiva
del error de cobertura y se mantiene dicha ventaja cuando se pasa de una
perturbación local a una estrategia dinámica con rebalanceo discreto?}
\end{quote}

Estas preguntas se evalúan separando precisión de pricing, precisión de Delta,
eficiencia muestral y desempeño de cobertura. En particular, no se presupone que
una mejora en la sensibilidad implique automáticamente una mejora uniforme del
precio ni del hedge dinámico.

\section{Objetivo general y objetivos específicos}
\label{sec:objetivos}

El objetivo general del TFM es:

\begin{quote}
\emph{evaluar si la incorporación de información diferencial mediante
Differential Machine Learning mejora, frente a un MLP convencional, la
aproximación del precio y especialmente de la Delta de una opción asiática
aritmética sobre capacidad de cómputo, y analizar en qué medida dicha mejora
se traduce en una reducción efectiva del riesgo de cobertura.}
\end{quote}

Para alcanzarlo se establecen los siguientes objetivos específicos:

\begin{enumerate}
    \item Caracterizar la serie observable de precios de alquiler H100 y evaluar
    si su comportamiento es compatible con una dinámica de reversión a la media.

    \item Construir un entorno reproducible de valoración de una opción asiática
    aritmética mediante Log-OU y Randomized Quasi-Monte Carlo.

    \item Generar y validar etiquetas de precio y Delta, contrastando la
    sensibilidad pathwise mediante diferencias finitas y controles numéricos
    independientes.

    \item Comparar de forma homogénea un MLP convencional y DML sobre tamaños de
    entrenamiento anidados y múltiples semillas, manteniendo separado el proceso
    de selección de la evaluación final.

    \item Analizar pricing, Delta, eficiencia muestral, errores extremos y
    sensibilidad alrededor de las fronteras de fixing.

    \item Evaluar la utilidad económica de las sensibilidades mediante una
    cobertura local y una estrategia dinámica con forward modelizado.
\end{enumerate}

\section{Contribuciones del trabajo}
\label{sec:contribuciones}

La aportación del TFM no consiste en proponer un nuevo modelo estocástico, un nuevo
algoritmo RQMC o una arquitectura neuronal original. Su contribución es integrar
herramientas consolidadas dentro de un experimento financiero controlado y
reproducible.

En primer lugar, se formula el precio por GPU-hora como una exposición económica
al coste de un servicio de cómputo y se distingue explícitamente entre la GPU
física, almacenable como bien de capital, y la GPU-hora, que representa capacidad
temporal. Sobre esta exposición se define una opción asiática aritmética como
instrumento hipotético de gestión del riesgo del coste medio.

En segundo lugar, se construye un entorno de pricing y sensibilidades que combina
una dinámica Log-OU, integración RQMC y una Delta pathwise validada mediante
diferencias finitas. El diseño separa los conjuntos utilizados para entrenamiento,
selección y evaluación y conserva una muestra de referencia de mayor precisión.

En tercer lugar, se realiza una comparación homogénea entre MLP y DML. Ambos
modelos comparten arquitectura, datos y protocolo de entrenamiento; la diferencia
sustantiva es el uso de supervisión diferencial. Esta estructura permite atribuir
con mayor claridad las diferencias observadas al tratamiento de la información de
Delta.

En cuarto lugar, el análisis no se limita a métricas de predicción. La calidad de
la Delta se traslada a dos pruebas de cobertura: una perturbación local, que mide
su capacidad para neutralizar pequeñas variaciones del subyacente, y un experimento
dinámico con rebalanceo discreto mediante un forward modelizado. De este modo se
distingue entre precisión estadística de una Greek y desempeño económico de un
sistema de cobertura.

Finalmente, el trabajo delimita expresamente su alcance. La serie H100 disponible
es corta, la medida de valoración es modelizada y el mercado de derivados sobre
compute es emergente. Por ello, las conclusiones se interpretan como evidencia
interna sobre el marco experimental propuesto y no como validación de precios
negociables observados en un mercado completo.

\section{Estructura del documento}
\label{sec:estructura_documento}

El Capítulo~2 revisa la literatura sobre activos no almacenables, opciones
asiáticas, reversión a la media, RQMC, redes neuronales para pricing, DML y
cobertura. El Capítulo~3 documenta las fuentes de datos, la selección del H100 y
las limitaciones de las referencias globales y chinas. El Capítulo~4 presenta la
metodología: contrato, Log-OU, generación de etiquetas, MLP, DML y experimentos
de cobertura. El Capítulo~5 analiza los resultados de pricing, Delta, eficiencia
muestral, errores extremos, fixings y cobertura. Finalmente, el Capítulo~6 responde
a las preguntas de investigación, resume las implicaciones financieras y expone
las principales limitaciones y extensiones.